\documentclass[12pt,openright,oneside,a4paper,brazil]{abntex2}
\usepackage[T1]{fontenc}
\usepackage[utf8]{inputenc}
\usepackage{lmodern}
\usepackage{microtype}
\usepackage{indentfirst}
\usepackage{setspace}
\usepackage{graphicx}
\usepackage{booktabs}
\usepackage{longtable}
\usepackage{hyperref}
\usepackage[alf]{abntex2cite}
\usepackage{geometry}
\geometry{a4paper,left=3cm,right=2cm,top=3cm,bottom=2cm}
\setlength{\parindent}{1.25cm}
\OnehalfSpacing
\hypersetup{hidelinks}

\titulo{Título provisório do projeto de pesquisa}
\autor{Nome do discente}
\local{Cidade}
\data{2026}
\instituicao{Instituição\\Curso}
\orientador{Nome do orientador}

\begin{document}
\imprimircapa
\imprimirfolhaderosto
\pdfbookmark[0]{\contentsname}{toc}
\tableofcontents*
\cleardoublepage

\chapter{Introdução}
\section{Problema de pesquisa}
Apresente o problema de forma delimitada e verificável.
\section{Objetivos}
\subsection{Objetivo geral}
\subsection{Objetivos específicos}
\section{Justificativa}

\chapter{Referencial teórico}
As afirmações devem usar citações verificadas e referências rastreáveis.

\chapter{Metodologia}
Descreva delineamento, participantes ou fontes, procedimentos, instrumentos, análise, ética, riscos e limitações.

\chapter{Resultados esperados}
Diferencie resultados esperados de resultados observados.

\chapter{Cronograma}
Use tabela conforme a norma tabular aplicável e registre a unidade de tempo.

\chapter{Considerações finais}
Retome o escopo sem afirmar conclusões além das evidências.

\postextual
\bibliography{referencias}

\end{document}